\documentclass{article}
\usepackage[utf8]{inputenc}
\usepackage{amsmath, amsthm, amssymb}
\usepackage{booktabs}
\usepackage{graphicx}
\usepackage{tcolorbox}
\usepackage{hyperref}
\tcbuselibrary{breakable}

\newtcolorbox{verification}{colback=green!5!white, colframe=green!40!black, title=Verification, breakable}

\title{ResLocoTransformer: Project Understanding Report}
\author{Gemini CLI Research Agent}
\date{\today}

\begin{document}

\maketitle

\section{Overview}
The goal of this project is to implement and train a transformer-based Reinforcement Learning (RL) policy for the Unitree Go2 quadruped robot using the MuJoCo physics engine. The project emphasizes visual-proprioceptive fusion and explores novel transformer architectures like the \texttt{LocoAttnResTransformer}.

\section{Environment: UnitreeMujocoGymEnv}
The environment is implemented in \texttt{envs/mujoco_env.py}.

\paragraph{Robot and Physics}
\begin{itemize}
    \item \textbf{Robot}: Unitree Go2.
    \item \textbf{Simulator}: MuJoCo 3.x.
    \item \textbf{Control}: Joint position targets (12-dim), converted to torques via PD control ($K_p=40, K_d=0.5$).
\end{itemize}

\paragraph{Observations}
The observation is a concatenated vector of size $84 + 16384 = 16468$.
\begin{itemize}
    \item \textbf{State (84-dim)}: Proprioceptive history (length 3).
    \begin{itemize}
        \item IMU: [roll, pitch, gyro_x, gyro_y] (4 dims $\times$ 3 steps = 12).
        \item Joint Angles: 12 joints $\times$ 3 steps = 36.
        \item Last Actions: 12 joints $\times$ 3 steps = 36.
    \end{itemize}
    \item \textbf{Depth (16384-dim)}: 4 stacked $64 \times 64$ depth frames.
    \begin{itemize}
        \item Rendered from a body-attached forward-facing camera.
        \item Log-transformed: $d' = \sqrt{\log(d+1)}$.
        \item Optionally z-scored.
    \end{itemize}
\end{itemize}

\paragraph{Reward Function}
The primary reward terms include:
\begin{itemize}
    \item \textbf{Velocity Tracking}: Exponential reward for matching target forward velocity.
    \item \textbf{Yaw Rate Tracking}: Exponential reward for matching target angular velocity.
    \item \textbf{Survival Bonus}: Constant positive reward per step.
    \item \textbf{Penalties}: Torque, action rate (smoothness), orientation (roll/pitch), height deviation, and lateral velocity/yaw rate.
\end{itemize}

\section{Model Architecture}
The architectures are defined in \texttt{models/nets.py} and \texttt{torchrl/networks/}.

\paragraph{LocoTransformerEncoder}
\begin{itemize}
    \item \textbf{Visual}: 3-layer CNN (NatureEncoder) $\to$ 16 tokens ($4 \times 4$ grid) of size \texttt{token_dim}.
    \item \textbf{State}: MLP $\to$ 1 token of size \texttt{token_dim}.
    \item Tokens are concatenated into a sequence of length 17.
\end{itemize}

\paragraph{LocoTransformer}
Standard Transformer layers followed by mean/max pooling of visual tokens and concatenation with the state token.

\paragraph{LocoAttnResTransformer}
Uses \textbf{Attention Residuals} (\texttt{AttnResLayer}). Standard additive residuals are replaced with depth-wise attention over all preceding layer outputs.
\[ x_{l+1} = \text{SoftmaxAttention}(q=f(x_l), K=H_l, V=H_l) \]
where $H_l = [x_0, x_1, \dots, x_l]$.

\section{Training Pipeline}
\begin{itemize}
    \item \textbf{Algorithm}: Proximal Policy Optimization (PPO).
    \item \textbf{Framework}: Custom \texttt{torchrl} library.
    \item \textbf{Normalization}: \texttt{NormObsWithImg} normalizes the state portion only.
    \item \textbf{Logging}: Weights and Biases (WandB) and local TensorBoard logs.
\end{itemize}

\section{Experiment Summary Table}
\begin{table}[h!]
\centering
\setlength{\tabcolsep}{8pt}
\renewcommand{\arraystretch}{1.2}
\begin{tabular}{lllrrc}
\toprule
ID & Date & Description & Commit & Metric & Status \\
\midrule
E000 & 2026-06-08 & Baseline Understanding & N/A & N/A & \textbf{Complete} \\
E001 & 2026-06-08 & Baseline Play Test & b8fc464 & N/A & \textbf{Failed (Arch Incompatible)} \\
\bottomrule
\end{tabular}
\caption{Summary of experiments.}
\end{table}

\section{Analysis of Baseline}
Attempting to run \texttt{scripts/play.py} on the baseline checkpoint (\texttt{00\_baseline}) resulted in a architecture mismatch. The current code assumes \texttt{LocoAttnResTransformer}, while the checkpoint seems to use a different (likely standard) transformer architecture.

\section{Experiment E002: Architecture Ablation (MLP vs. Visual Pool vs. AttnRes)}

\paragraph{Goal}
Decide whether visual tokens and Attention Residuals (AttnRes) are worth their
extra compute on rough terrain, by isolating the contribution of (a) visual
input at all and (b) attention over visual tokens vs. simple pooling.

\paragraph{Hypothesis}
Visual terrain information should help on \texttt{scene\_rough.xml} beyond
what proprioception alone (MLP arm) provides, since the policy cannot see
upcoming terrain changes otherwise. Whether attention over visual tokens
(AttnRes) beats simple mean-pooling of the same tokens (Visual Pool) tests
whether cross-token interaction is doing useful work, or whether the visual
encoder's per-token features are already sufficient once pooled.

\paragraph{Method}
Three arms, identical in reward weights, \texttt{use\_*} flags, environment
settings (\texttt{scene\_rough.xml}), PPO hyperparameters, and epoch budget
(2000 epochs); only the policy architecture differs:
\begin{itemize}
    \item \textbf{MLP} (\texttt{configs/ablation\_mlp.json}): state-only baseline, \texttt{policy\_type: "mlp"}, \texttt{get\_image: false}. 222k policy parameters.
    \item \textbf{Visual Pool} (\texttt{configs/ablation\_visual\_pool.json}): visual encoder produces tokens, mean-pooled (\texttt{transformer\_params: []}, \texttt{max\_pool: false}), no attention layers. 288k parameters.
    \item \textbf{AttnRes} (\texttt{configs/ablation\_attnres.json}): same visual encoder, plus 2 \texttt{AttnResLayer} layers (\texttt{[4, 128]} each, \texttt{attn\_res\_heads: 4}). 422k parameters.
\end{itemize}
Each arm is run with 3 seeds (0, 1, 2) via \texttt{scripts/run\_ablations.py}
(\texttt{--seeds 0 1 2 --vec-env-nums 8 --proc-nums 8}, reduced from the
script's default of 16 to fit the 8GB VRAM budget of the local RTX 3070 --
this parallelism setting is held identical across all 9 runs so it cannot
confound the architecture comparison).

\paragraph{Implementation}
\texttt{scripts/run\_ablations.py} (new), \texttt{configs/ablation\_\{mlp,visual\_pool,attnres\}.json}.

\paragraph{Results}
\textit{Sweep in progress on the RTX 3070 (9 runs: 3 arms $\times$ 3 seeds,
2000 epochs each). Table to be filled in on completion.}

\paragraph{Analysis}
TBD -- pending sweep completion. Per project convention, reward alone will
not decide a winner: \texttt{reward/foot\_slip}, \texttt{reward/foot\_clearance},
\texttt{reward/feet\_air\_time}, and \texttt{diag/base\_height} (converged,
last 10\% of epochs) will be checked alongside reward, and the \texttt{best}
checkpoint of each arm will be replayed with \texttt{scripts/play.py} before
calling a winner, since reward is currently known to be an unreliable signal
in this codebase.

\paragraph{Next steps}
Run \texttt{scripts/compare\_experiments.py --log-dir ./log\_ablation} once
all 9 runs finish; replay best checkpoints; fill in results table and analysis
above.

\begin{verification}
\textbf{What:} All three ablation configs build and run a forward/backward
pass without crashing, for at least 2 epochs each.

\textbf{Method:} Tier-1 smoke test -- ran each config for 2--3 epochs with
\texttt{vec\_env\_nums=4} before launching the full sweep.

\textbf{Script:} manual invocation of \texttt{scripts/train.py} with a
temporary reduced-epoch copy of each config (not checked in).

\textbf{Outcome:} pass; all three arms complete epochs without traceback.
The tensorboardX ``NaN or Inf found in input tensor'' warning at epoch 0 is
benign (comes from \texttt{Running\_Training\_Average\_Rewards} having no
data yet on the very first epoch, not from the policy/value networks).
\end{verification}

\section{Experiment Summary Table (Updated)}
\begin{table}[h!]
\centering
\setlength{\tabcolsep}{8pt}
\renewcommand{\arraystretch}{1.2}
\begin{tabular}{lllrrc}
\toprule
ID & Date & Description & Commit & Metric & Status \\
\midrule
E000 & 2026-06-08 & Baseline Understanding & N/A & N/A & \textbf{Complete} \\
E001 & 2026-06-08 & Baseline Play Test & b8fc464 & N/A & \textbf{Failed (Arch Incompatible)} \\
E002 & 2026-08-12 & Architecture Ablation (MLP/VisualPool/AttnRes) & d490d18 & Running\_Average\_Rewards + gait diagnostics & \textbf{Running} \\
\bottomrule
\end{tabular}
\caption{Summary of experiments.}
\end{table}

\end{document}
